% Глава: 23-docker-resources
% Статус: наполнено (первый прогон).
\chapter{Ресурсы контейнеров}
\label{ch:23-docker-resources}

\section{Изоляция ресурсов: cgroups}

Контейнеры делят ядро хоста; чтобы один «не съел»
всё, Linux использует \textbf{cgroups} (control
groups) --- механизм ядра, ограничивающий ресурсы
(память, CPU, IO) для группы процессов. Docker
управляет cgroups за вас через флаги.

Лимиты задают долю ресурсов вашей ВМ на \fpath{pglab.nsu.ru},
чтобы один контейнер не забрал всё у соседних.

\section{Лимиты CPU}

\begin{minted}[bgcolor=gray!6]{bash}
docker run --cpus 1.5 myimage          # ограничить 1.5 ядрами
docker run --cpuset-cpus 0,1 myimage   # закрепить за конкретными CPU
\end{minted}

\cmd{--cpus} задаёт долю (дробное число допустимо),
\cmd{--cpuset-cpus} --- список конкретных ядер.

\section{Лимиты памяти}

\begin{minted}[bgcolor=gray!6]{bash}
docker run --memory 512m myimage          # максимум 512 МиБ
docker run --memory 512m --memory-swap 1g myimage  # + подкачка (swap)
\end{minted}

Если процесс превысит лимит памяти, ядро сработает
\textbf{OOM-killer} и завершит процесс с ошибкой
(\cmd{Killed}). Проверьте логи:
\cmd{docker inspect} покажет \cmd{OOMKilled: true}.

\section{Мониторинг}

\begin{minted}[bgcolor=gray!6]{bash}
docker stats             # живая статистика CPU и памяти по контейнерам
docker stats --no-stream # однократно
docker top <id>          # процессы внутри контейнера
docker inspect <id>      # полная конфигурация и лимиты
\end{minted}

\section{Типовая схема ограничений}

\begin{figure}[htb]
  \centering
  \begin{tikzpicture}[node distance=1.3cm]
    \node[layer, minimum width=7cm, minimum height=2.8cm] (host) {};
    \node[font=\small] at (host) {хост: CPU + RAM};
    \node[process, minimum width=1.6cm] (c1) at ([yshift=-0.4cm]host.west) [xshift=0.9cm] {app 1};
    \node[process, right=0.3cm of c1, minimum width=1.6cm] (c2) {app 2};
    \node[font=\tiny, below=1pt of c1] {cpus=1.0, mem=512m};
    \node[font=\tiny, below=1pt of c2] {cpus=0.5, mem=256m};
  \end{tikzpicture}
  \caption{Каждый контейнер имеет лимиты CPU и памяти (cgroups).}
  \label{fig:cgroups}
\end{figure}

\section{Типичные ошибки}
\begin{itemize}
  \item Контейнер убивается «без причины» --- смотрите
    \cmd{OOMKilled}; увеличьте \cmd{--memory} или
    оптимизируйте память приложения;
  \item Лимит CPU слишком низкий --- приложение
    «замедляется»; измерьте \cmd{docker stats};
  \item Не заданы лимиты --- один контейнер может
    занять всю память хоста;
  \item \cmd{--memory-swap} без \cmd{--memory} ---
    некорректная конфигурация.
\end{itemize}

\section{Связь с остальным курсом}

Лимиты настраиваются и в \cmd{docker-compose}
(\nameref{ch:22-docker-compose}); мониторинг процессов
--- \nameref{ch:12-processes}; распределение памяти ---
\nameref{ch:16-virtual-memory}.

\section{Задачи для самопроверки}
\begin{enumerate}
  \item Что такое cgroups и зачем они нужны
    контейнерам?
  \item Как ограничить контейнер 1 CPU и 512 МБ
    памяти?
  \item Как узнать, был ли контейнер убит
    OOM-killer'ом?
  \item Чем \cmd{--cpus} отличается от
    \cmd{--cpuset-cpus}?
  \item Как посмотреть использование ресурсов
    контейнеров в реальном времени?
\end{enumerate}
